\begin{abstracts}        

    
Artificial intelligence is shaping educational practice, while policy frameworks seek
to guide its responsible adoption across learning environments. This research
investigates how computational text analysis can identify differences between
AI-in-education policy and public sentiment. A multilingual text-as-data framework
was developed using policy and sentiment evidence from Ireland, France, and European
sources. The methodology combines topic modelling, semantic causal-claim coverage,
and causal-frame network analysis, with synthetic sentiment used for robustness
testing. An exploratory LLM-assisted stage operates on the sentence inventory to
provide evidence-grounded interpretation of causal and conditional relationships.
Together, these methods examine whether comparable causal claims occur across policy
and sentiment and whether educational concerns are organised through causes,
requirements, mechanisms, risks, and outcomes. Results show semantic correspondence
between policy and sentiment, alongside selective differences in causal-claim
coverage and causal-frame organisation. The differences concern learning and
assessment, curriculum and oversight, and causal roles assigned to institutional
mechanisms and consequences. Policy places emphasis on enabling mechanisms,
requirements, and intended improvements, whereas sentiment gives greater weight to
expressed causes and consequences. The LLM-assisted stage shows that less formulaic
causal or conditional relationships remain outside fixed cue extraction, while
providing interpretations of evidence relationships represented by semantic
similarity. These findings do not establish real-world causal effects, policy
effectiveness, or across-the-board LLM superiority. Instead, they demonstrate the
value of combining reproducible semantic and structural causal NLP measurements with
traceable, evidence-grounded interpretation for analysing one observable dimension of
the policy--practice gap in AI-in-education governance.

\end{abstracts}
